\section{V5 Hardening and Profile Audit}
\label{sec:v5-hardening}

V5 adds independent checks around the frozen V4 experiment without silently changing its primary classifications. The hardening layer is additive: it records chronology, source and destination fault sensitivity, holdout membership, profile settings, and implementation disagreement so that a reviewer can distinguish stronger evidence from a changed result denominator.

\subsection{Profile interpretation and condition comparison}
\label{sec:v5-profile}

The V4 LibreOffice condition used \texttt{UseTaggedPDF=true}, \texttt{PDFUACompliance=true}, and \texttt{SelectPdfVersion=1}. The V5 profile audit documents that the latter selects PDF/A-1b rather than an ordinary default PDF export. V5 condition B repeated all 65 LibreOffice fixtures with \texttt{SelectPdfVersion=0}, while retaining the tagging and PDF/UA settings. All 65 condition-B conversions produced valid outputs; all were tagged, and the raw PDF/A marker count was zero. The machine-derived classifier produced zero semantic classification changes relative to V4. This is evidence of no material effect on the tested semantic classifier, not evidence of byte identity, complete PDF/UA conformance, or equivalence of every PDF serialization detail.

\subsection{Holdout accounting and fault sensitivity}
\label{sec:v5-holdout}

The V5 holdout audit derives fixture membership from the frozen manifests rather than from prose counts. The 130 primary fixture--pipeline cases contain 52 fixture-level holdout cases (26 per primary pipeline) and 78 non-holdout cases (39 per pipeline). The holdout is not pooled with calibration controls, and the supplementary Word arm remains outside the primary denominator.

The source mutation audit detected 39 of 39 deliberate source mutations. Destination-side mutation testing applied the family-specific operators across all 65 atomic LibreOffice PDFs, spanning five variants per family and multiple structure/value locations. Fifty-nine attempts produced valid destination mutants; 46 changed the frozen classification, 13 survived without a classification change, six attempts were unsupported, and two valid mutants produced measurement errors. Sixty-five no-op controls retained their baseline classifications with zero false-positive controls. These tests establish sensitivity and expose blind spots for the recorded fault models, not exhaustive coverage of all converter or PDF-parser faults.

\subsection{Independent diagnostic cross-check and calibration}
\label{sec:v5-validation}

An independent lightweight raw/PyMuPDF implementation inspected 26 rows, one per feature family per primary pipeline. It agreed with the original classification on 8 rows and disagreed on 18. The disagreement is informative: the second implementation intentionally lacked some exact-value and association logic, so it frequently abstained or could not establish the richer contract. It is therefore reported as an independent diagnostic cross-check, not as an external adjudicated ground truth and not as a replacement for the frozen comparator.

The 78-case calibration audit produced 60 correct decisions, zero false positives, five conservative false negatives, and 13 correct abstentions. The calibration result supports the claimed conservative behavior for the tested controls; it is not a sensitivity, specificity, or population-performance estimate.

\subsection{What V5 establishes}
\label{sec:v5-interpretation}

Taken together, the V5 checks strengthen four claims. First, the primary accounting is prospectively separable into holdout and non-holdout cases. Second, the comparator detects the tested source and destination perturbations. Third, the main result is not an artifact of the documented LibreOffice PDF/A profile interpretation. Fourth, a second technical implementation exposes where exact cross-format equivalence is difficult rather than hiding that uncertainty.

The hardening does not convert the study into a global accessibility benchmark. The atomic corpus remains controlled and purposive; the realistic documents are project-created; Google Docs remains a cloud condition whose backend cannot be pinned; the Word arm is supplementary; reading order remains outside the tested attribute set; and no disabled-user or screen-reader task study was performed. These are properties of the research question and evidence boundary, not packaging defects. The final claim is consequently precise: A11yPreserve provides an auditable source-contract testing methodology for selected machine-verifiable accessibility semantics across the recorded DOCX-to-PDF conditions.
